\begin{table*}[t]
\caption{\label{tab:III}Ionic-relaxation benchmark for the tetramer $P$-substitution synergy grid (four fixed-cell geometry-optimization corners: pristine and $P$ at $\epsilon{=}0$ and $+3$\%, same biaxial cell).
This benchmark is \emph{sign-qualitative only}: the rigid reference uses legacy PBE without D3, whereas relaxed corners use PBE+D3, so $|\mathcal{S}_{\mathrm{relaxed}}|/|\mathcal{S}_{\mathrm{rigid}}|$ is \emph{not} a quantitative retention ratio until a matched-functional rigid grid is recomputed. We test whether $\mathcal{S}$ retains \emph{sign} \emph{within the tetramer grid only}; it does \emph{not} calibrate periodic main-text $\mathcal{S}(n)$, which uses different dopant concentration and boundary conditions.
Rigid reference energies: archived fixed-coordinate tetramer single-point grid (legacy PBE, no D3). Relaxed energies: PBE+D3 fixed-cell geometry optimization (400~Ry) from archived inputs in the open repository.
$\mathcal{S}$ follows Eq.~\eqref{eq:synergy_order} with $N{=}240$ atoms.}
\begin{ruledtabular}
\begin{tabular}{lcc}
Quantity & Rigid (legacy tetramer) & Relaxed (geometry opt.) \\
\hline
$E$(pristine, $\epsilon{=}0$) ($E_{\mathrm{h}}$) & $-1363.5512$ & $-1363.9551$ \\
$E$(pristine, $\epsilon{=}3$\%) ($E_{\mathrm{h}}$) & $-1363.5522$ & $-1363.9550$ \\
$E$(P, $\epsilon{=}0$) ($E_{\mathrm{h}}$) & $-1370.4965$ & $-1374.4415$ \\
$E$(P, $\epsilon{=}3$\%) ($E_{\mathrm{h}}$) & $-1370.4890$ & pending \\
$\mathcal{S}$ (meV/atom) & $+0.96$ & pending \\
Sign of $\mathcal{S}$ vs.\ rigid & \multicolumn{2}{c}{pending (3/4 geometry opt.; $P$ at $\epsilon{=}3$\% in progress)} \\
\end{tabular}
\end{ruledtabular}
\end{table*}
